% SECTION I: INTRODUCTION
% Polished with crisp contributions

\section{Introduction}
\label{sec:intro}

\IEEEPARstart{R}{econfigurable} intelligent surfaces (RIS) have emerged as a transformative technology for sixth-generation (6G) wireless networks, offering unprecedented control over the electromagnetic propagation environment through passive beamforming~\cite{Basar2019Access}. By dynamically adjusting the phase shifts of numerous low-cost reflecting elements, RIS can enhance coverage, improve spectral efficiency, and enable new applications across millimeter-wave and terahertz bands~\cite{Wu2019TWC}.

\subsection{The Latency Blind Spot in RIS Optimization}

The vast majority of RIS optimization research focuses on maximizing performance under the implicit assumption that computation time is negligible or exogenous to algorithm design. Classical approaches include alternating optimization (AO)~\cite{Wu2019TWC}, weighted minimum mean-square error (WMMSE)~\cite{Guo2020TCOM}, and projected gradient descent (PGD). Recent work has explored deep learning for direct phase prediction~\cite{Huang2020TWC,Elbir2020WCL}, motivated by the potential to bypass iterative optimization entirely.

However, these studies share a critical limitation: \textbf{they evaluate methods by convergence quality alone}, measuring performance at the endpoint of optimization without considering the time required to reach it. In static channels this evaluation is appropriate. In mobile scenarios it is fundamentally flawed.

\subsection{The Mobility-Latency Paradox}

Mobile users experience time-varying channels with finite coherence time $T_c$. At carrier frequency $f_c = 3$ GHz and pedestrian velocity $v = 1.5$ m/s, coherence time is approximately 28.2 ms~\cite{rappaport_wireless}. Any computation exceeding this duration renders the optimized phase outdated before deployment.

This creates a fundamental trade-off: \textit{converged solutions require many iterations, but by the time convergence is achieved, the channel has changed}. Classical projected gradient ascent achieves near-optimal performance with 1000 iterations but requires approximately 5600 ms---199$\times$ longer than coherence time. By the time the optimal phase is computed, it is applied to a completely different channel state.

The consequence is striking: \textbf{designing for convergence systematically misallocates computational resources}. Methods are ranked by their quality at $t=0$ (when optimization begins) rather than their quality at $t=\tau$ (when results become available), where $\tau$ is the algorithm-dependent processing delay.

\subsection{Gap in Existing Literature}

While mobility and channel aging have been studied extensively in MIMO systems~\cite{Jindal2006IT}, the interaction between \textit{computation delay} and \textit{channel aging} remains largely unexplored in the RIS literature. Existing mobility-aware RIS work~\cite{Huang2021TWC,Papazafeiropoulos2022TCOM} typically assumes zero or negligible computation time, focusing instead on tracking accuracy or channel prediction horizons. Predictive beamforming schemes~\cite{Ma2020WCL,Liu2022TCOM,Zhao2023WCL} forecast future channels $\hat{\mbf{H}}(t + T_{\text{proc}})$ but \textit{fail to eliminate staleness} because prediction itself consumes time---the cumulative delay $T_{\text{pred}} + T_{\text{opt}}$ still renders the phase configuration stale.

To our knowledge, no prior work has:
\begin{itemize}
    \item \textbf{Quantified} the value of fresh beamforming relative to stale optimized phase in RIS systems
    \item \textbf{Proven} the existence of an intrinsic optimum (finite iteration count that maximizes effective rate)
    \item \textbf{Demonstrated} that method ranking systematically reverses under equal-millisecond comparison
    \item \textbf{Provided} reproducible measurements with documented experimental scripts and sabotage-tested audit framework
\end{itemize}

\subsection{Contributions}

This paper fills this gap through rigorous latency-aware analysis of RIS phase optimization in mobile multi-user scenarios. Our main contributions are:

\begin{enumerate}
    \item \textbf{We prove fresh beamforming dominates stale optimized phase} (Lemma~\ref{lem:fresh_dominance}): Fresh beamforming with instantaneous CSI is 11.7--14.2$\times$ more valuable than stale converged phase at coherence time $T_c = 28.2$ ms. This corrects a previously reported but non-reproducible claim of ``17$\times$ gain'' for which no executable script existed. Our measurement is fully reproducible via scripts/mumiso\_beam\_staleness.py with 10{,}000 independent channel realizations.

    \item \textbf{We prove the existence of an intrinsic optimum at finite iteration count} (Theorem~\ref{thm:intrinsic_optimum}): We identify $L^* \approx 8$ gradient steps as the intrinsic optimum for pedestrian scenarios ($v = 1.5$ m/s, $f_c = 3$ GHz), achieving 5.380 bps/Hz effective rate. Designing for deep convergence ($L=48$ steps) costs 5.9\% in achievable rate despite superior t=0 performance, proving that \textit{mathematical convergence is not physically optimal} under mobility.

    \item \textbf{We prove method ranking systematically reverses under equal-millisecond constraints} (Proposition~\ref{prop:ranking_reversal}): Per-step latency variation (2.8$\times$ between learned updates and Adam optimizer) causes \textit{complete ranking reversal}. Under equal-iteration comparison, learned gradient updates win at all five budgets ($L=1,2,4,8,16$) with 68.3--83.7\% win rates. Under equal-millisecond comparison, all five learned configurations are dominated by classical Adam optimizer across five independent trials, confirming that per-iteration latency is a first-class design constraint.

    \item \textbf{We provide a reproducibility framework with sabotage-tested integrity}: All 41 audited claims are backed by executable scripts in a version-controlled repository. Sabotage testing (writing untrained weights to trained checkpoints) caused 7 out of 41 assertions to fail, confirming the audit measures genuine model quality rather than cached constants. We document measurement traps (machine load contamination, timing jitter, prefix non-stability) that affect latency-dependent conclusions, and provide idle-machine validation protocols to prevent contamination.
\end{enumerate}

\subsection{Paper Organization}

Section~\ref{sec:related_work} positions our work relative to classical optimization, deep learning, mobility-aware design, and predictive beamforming literature. Section~\ref{sec:system} presents the 3GPP-compliant system model, Jakes mobility framework, and latency-aware effective rate formulation. Section~\ref{sec:analysis} establishes the theoretical foundation through three main results (Lemma~\ref{lem:fresh_dominance}, Theorem~\ref{thm:intrinsic_optimum}, Proposition~\ref{prop:ranking_reversal}). Section~\ref{sec:measurements} describes the reproducibility framework and measurement protocols. Section~\ref{sec:results} validates the theoretical findings through extensive simulations. Section~\ref{sec:conclusion} concludes the paper and discusses future directions.
